\subsection{Regime 2: dilute epitopes (no antibody depletion)}

As in Regime 1, let $S_c=A+B$ and $S_t=A'+B$ be the total target-epitope
concentrations in the control and treatment conditions, and let
$S\in\{S_c,S_t\}$ denote either one. In Regime 1, antibody was the scarce
reagent: nearly every antibody molecule found a target, so the amount of
target captured was fixed at about $Y$ and the \emph{fraction} of target
captured depended on $S$. Here we consider the opposite situation: the
epitopes are so dilute that binding them removes only a negligible amount
of antibody from solution. The free antibody concentration is then set
by the experimenter ($Y_f\approx Y$) rather than by the chromatin, and,
as we will see, every capture efficiency becomes independent of $S$.

We assume the following for both conditions (i.e., for each
$S\in\{S_c,S_t\}$).

\begin{enumerate}
    \item Target epitopes are dilute relative to $K_{d1}$:
    $$
    S\ll K_{d1}.
    $$
    To interpret the ratio $S/K_{d1}$, rearrange the mass-action relation
    $K_{d1}=Y_f(A_f+B_f)/(Y_A+Y_B)$ to give the ratio of target-bound to
    free antibody:
    $$
    \frac{Y_A+Y_B}{Y_f}=\frac{A_f+B_f}{K_{d1}}.
    $$
    When the target is unsaturated, nearly all target is free
    ($A_f+B_f\approx S$), so $S/K_{d1}$ is approximately the ratio of
    target-bound antibody to free antibody. The assumption therefore says
    that for every antibody molecule sitting on a target epitope, many more
    are free in solution. (If the target is partly saturated, fewer target
    epitopes are free, $A_f+B_f<S$, so $S/K_{d1}$ is an upper bound on this
    ratio.) Note the contrast with Regime 1, which required the
    opposite inequality, $K_{d1}\ll S-Y$.

    \item Off-target epitopes are dilute relative to $K_{d2}$:
    $$
    C+D\ll K_{d2}.
    $$
    By the same logic applied to the off-target mass-action relation,
    $$
    \frac{Y_C+Y_D}{Y_f}=\frac{C_f+D_f}{K_{d2}}\le\frac{C+D}{K_{d2}},
    $$
    so $(C+D)/K_{d2}$ bounds the ratio of off-target-bound to free
    antibody. Because $K_{d2}\gg K_{d1}$, this assumption is usually much
    easier to satisfy than assumption 1: the off-target pool may be far
    larger than the target pool and still bind little antibody.
\end{enumerate}
In summary, the two assumptions say that the two places antibody can go
other than solution---target and off-target epitopes---each hold a small
amount of antibody compared with the amount that stays free. Note that
neither assumption involves $Y$: in this regime, the antibody may be
limiting, intermediate, or in excess relative to the dissociation
constants; it is only required to be in excess of what the epitopes can
soak up. Because $S_t<S_c$ when $A'<A$, the control condition is the more
demanding one, so it suffices to check the assumptions with $S=S_c$.

We now derive $Y_f$, $\theta_1$, and $\rho$ for either condition. (For the
treatment condition, $\theta_1$, $\theta_2$, $Y_f$, and $\rho$ should be
read as their primed counterparts.)

\textbf{Antibody conservation.} As derived in Regime 1, substituting
$Y_A+Y_B=S\theta_1$ and $Y_C+Y_D=(C+D)\theta_2$ into
$Y=Y_f+Y_A+Y_B+Y_C+Y_D$ gives
\begin{equation}
Y=Y_f+S\theta_1+(C+D)\theta_2.
\label{eq:r2-1}
\end{equation}

\textbf{Step 1: bound the antibody on epitopes.} Dropping $Y_f$ from the
denominator of $\theta_1=Y_f/(Y_f+K_{d1})$ makes the denominator smaller
and hence the fraction larger:
$$
\theta_1=\frac{Y_f}{Y_f+K_{d1}}\le\frac{Y_f}{K_{d1}}.
$$
Multiplying both sides by $S$,
$$
S\theta_1\le S\frac{Y_f}{K_{d1}}=\frac{S}{K_{d1}}\,Y_f.
$$
By the same logic with $\theta_2=Y_f/(Y_f+K_{d2})\le Y_f/K_{d2}$,
$$
(C+D)\theta_2\le\frac{C+D}{K_{d2}}\,Y_f.
$$

\textbf{Step 2: $Y_f\approx Y$.} Define the relative amount of
epitope-bound antibody
$$
\delta=\frac{S\theta_1+(C+D)\theta_2}{Y_f},
$$
so that \eqref{eq:r2-1} reads $Y=Y_f+\delta Y_f=Y_f(1+\delta)$. Adding the
two inequalities from Step 1 and dividing by $Y_f$,
\begin{equation}
0\le\delta
\le\frac{S}{K_{d1}}+\frac{C+D}{K_{d2}}
\equiv\bar\delta
\ll1,
\label{eq:r2-2}
\end{equation}
where $\delta\ge0$ because every term is nonnegative, and the final
$\ll$ follows from assumptions 1 and 2 (a sum of two small numbers is
small). Solving $Y=Y_f(1+\delta)$ for $Y_f$ and using
$1\le1+\delta\le1+\bar\delta$,
\begin{equation}
\frac{Y}{1+\bar\delta}\le Y_f=\frac{Y}{1+\delta}\le Y.
\label{eq:r2-3}
\end{equation}
Since $\bar\delta\ll1$, $Y_f\approx Y$: the free antibody is essentially
all of the antibody. (Equivalently, the fraction of antibody that is
bound, $(Y-Y_f)/Y=\delta/(1+\delta)$, is at most $\bar\delta$.)

\textbf{Step 3: capture efficiencies.} Substituting $Y_f\approx Y$ into
the expressions for the bound fractions,
$$
\theta_1\approx\frac{Y}{Y+K_{d1}},
\qquad
\theta_2\approx\frac{Y}{Y+K_{d2}}.
$$
Neither depends on $S$. In particular, the fraction of target epitope
captured is the same in both conditions, so the IP yield \emph{per
epitope} does not change when the treatment removes target. (The total
target captured, $S\theta_1$, is proportional to $S$ and so does
decrease.) This is the key difference from Regime 1, where the total
captured was fixed at about $Y$ and the per-epitope efficiency
$\theta_1\approx Y/S$ rose when $S$ fell.

\textbf{Step 4: $\rho$.} Substituting $Y_f\approx Y$ into
$\rho=(Y_f+K_{d1})/(Y_f+K_{d2})$ gives
\begin{equation}
\rho\approx\rho'\approx\bar\rho\equiv\frac{Y+K_{d1}}{Y+K_{d2}}.
\label{eq:r2-4}
\end{equation}
Since $\bar\rho$ depends only on $Y$, $K_{d1}$, and $K_{d2}$---not on
$S$---it is the same for the control and treatment conditions,
regardless of $A'$.

To make ``$\approx$'' precise, we bound $\rho$ above and below. First,
$\rho$ increases with $Y_f$: by the quotient rule,
$$
\frac{d\rho}{dY_f}
=\frac{(Y_f+K_{d2})-(Y_f+K_{d1})}{(Y_f+K_{d2})^2}
=\frac{K_{d2}-K_{d1}}{(Y_f+K_{d2})^2}
>0,
$$
since $K_{d2}>K_{d1}$. Because $Y_f\le Y$ by \eqref{eq:r2-3}, it follows
that $\rho\le\bar\rho$. For the lower bound, substitute the lower bound
$Y_f\ge Y/(1+\bar\delta)$ from \eqref{eq:r2-3} and multiply the numerator
and denominator by $1+\bar\delta$:
$$
\rho
\ge\frac{\dfrac{Y}{1+\bar\delta}+K_{d1}}{\dfrac{Y}{1+\bar\delta}+K_{d2}}
=\frac{Y+(1+\bar\delta)K_{d1}}{Y+(1+\bar\delta)K_{d2}}.
$$
The numerator is at least $Y+K_{d1}$, and the denominator is at most
$(1+\bar\delta)Y+(1+\bar\delta)K_{d2}=(1+\bar\delta)(Y+K_{d2})$, so
$$
\rho\ge\frac{Y+K_{d1}}{(1+\bar\delta)(Y+K_{d2})}=\frac{\bar\rho}{1+\bar\delta}.
$$
Together,
\begin{equation}
\frac{\bar\rho}{1+\bar\delta}\le\rho\le\bar\rho,
\label{eq:r2-5}
\end{equation}
so $\rho$ agrees with $\bar\rho$ to within a relative error of at most
$\bar\delta\ll1$. The same holds for $\rho'$ (with $\bar\delta$ computed
using $S_t$, which is smaller than with $S_c$).

\textbf{Direction of the (small) difference.} Although $\rho'\approx\rho$,
the exact values still move in the same direction as in Regime 1. Write
the right-hand side of \eqref{eq:r2-1} as a function of $Y_f$ and $S$:
$$
g(Y_f,S)=Y_f+S\frac{Y_f}{Y_f+K_{d1}}+(C+D)\frac{Y_f}{Y_f+K_{d2}}.
$$
Each term increases with $Y_f$ (for instance,
$\frac{d}{dY_f}\frac{Y_f}{Y_f+K}=\frac{K}{(Y_f+K)^2}>0$), and $g$ increases
with $S$. If $A'<A$, then $S_t<S_c$, so
$$
g(Y_f,S_t)<g(Y_f,S_c)=Y=g(Y_f',S_t).
$$
Since $g(\cdot,S_t)$ is increasing, this requires $Y_f<Y_f'$: with less
target to bind, slightly more antibody stays free. Because $\rho$
increases with $Y_f$, $\rho<\rho'$. Combining with \eqref{eq:r2-5} for
both conditions,
\begin{equation}
\rho<\rho'\le\bar\rho\le(1+\bar\delta)\rho,
\label{eq:r2-6}
\end{equation}
where $\bar\delta$ is evaluated for the control ($S=S_c$). Thus $\rho'$
exceeds $\rho$, but by a relative amount of at most $\bar\delta\ll1$.

\textbf{Sub-limits of $\bar\rho$.} Factoring $K_{d1}$ out of the numerator
and $K_{d2}$ out of the denominator of \eqref{eq:r2-4},
$$
\bar\rho
=\frac{K_{d1}(1+Y/K_{d1})}{K_{d2}(1+Y/K_{d2})}
=r\,\frac{1+Y/K_{d1}}{1+Y/K_{d2}}.
$$
\begin{itemize}
    \item $Y\ll K_{d1}$ (antibody below the target $K_d$): both
    $Y/K_{d1}\ll1$ and $Y/K_{d2}\ll1$, so $\bar\rho\approx r$. Both
    epitope classes are unsaturated, $\theta_1\approx Y/K_{d1}$ and
    $\theta_2\approx Y/K_{d2}$, and their ratio is set purely by the
    affinities.
    \item $K_{d1}\ll Y\ll K_{d2}$: in the numerator,
    $Y+K_{d1}\approx Y$; in the denominator, $Y+K_{d2}\approx K_{d2}$. So
    $\bar\rho\approx Y/K_{d2}$. The target is saturated
    ($\theta_1\approx1$), while off-target binding is still in its linear
    regime ($\theta_2\approx Y/K_{d2}$), so off-target capture grows in
    proportion to $Y$ while target capture cannot grow further.
    \item $Y\gg K_{d2}$: both $Y+K_{d1}\approx Y$ and $Y+K_{d2}\approx Y$,
    so $\bar\rho\approx1$. Both classes are saturated and the antibody no
    longer discriminates between them.
\end{itemize}
As $Y$ increases, $\bar\rho$ rises monotonically from $r$ to 1 (by the
same derivative computed above, with $Y$ in place of $Y_f$).

\subsubsection{ChIP normalization}

The read-count expressions derived in Regime 1 (expected reads $n_i$ at
target locus $i$ and $N_s$ on the spike-in genome) did not use any
regime-specific assumptions, so they apply here unchanged.

\textbf{Scaling by spike-in reads (ChIP-Rx).} From Regime 1,
$$
\frac{n_i'/N_s'}{n_i/N_s}
=\frac{a_i'}{a_i}\cdot\frac{B+\rho D}{B+\rho' D}.
$$
Substituting $\rho\approx\rho'\approx\bar\rho$ from \eqref{eq:r2-4}, the
numerator and denominator of the bias factor are both approximately
$B+\bar\rho D$, so
$$
\frac{B+\rho D}{B+\rho' D}\approx\frac{B+\bar\rho D}{B+\bar\rho D}=1,
\qquad
\frac{n_i'/N_s'}{n_i/N_s}\approx\frac{a_i'}{a_i}.
$$
More precisely, by \eqref{eq:r2-6}, $\rho D<\rho' D\le(1+\bar\delta)\rho D$.
Adding $B$ and using $B\le(1+\bar\delta)B$,
$$
B+\rho D<B+\rho' D\le(1+\bar\delta)(B+\rho D),
$$
and taking reciprocals and multiplying by $B+\rho D$,
$$
\frac{1}{1+\bar\delta}\le\frac{B+\rho D}{B+\rho' D}<1.
$$
So spike-in normalization recovers the true fold change to within a
relative error of at most $\bar\delta$, \emph{even if off-target spike-in
reads are not negligible} ($\rho D$ comparable to or larger than $B$).
This contrasts with Regime 1, where the bias factor equals 1 only if
$\rho D\ll B$. The reason is that in this regime the off-target spike-in
reads are captured with the same efficiency in both conditions, so they
scale the spike-in read count by the same factor in both libraries and
cancel in the ratio. (The direction of the residual bias, $<1$, is the
same as in Regime 1, but its size is bounded by $\bar\delta$.) Scaling by
spike-in target reads $N_{s,B}$ also works, as it did in Regime 1.

\textbf{Scaling by total reads (RPM).} From Regime 1, substituting
$\rho'\approx\rho$,
$$
\frac{n_i'/N'}{n_i/N}
=\frac{a_i'}{a_i}\cdot\frac{A+B+\rho(C+D)}{A'+B+\rho'(C+D)}
\approx\frac{a_i'}{a_i}\cdot\frac{A+B+\rho(C+D)}{A'+B+\rho(C+D)}.
$$
If $A'<A$, the numerator of the second factor exceeds its denominator, so
the factor is greater than 1 and partially cancels the global decrease
(fully, $\approx A/A'$, if the spike-in and off-target terms are small).
RPM normalization still hides global changes in this regime.

\textbf{The cost of this regime.} Dilute epitopes make normalization
robust, but they come at a price.
\begin{itemize}
    \item \emph{Low yield when $Y\ll K_{d1}$.} The fraction of target
    captured is $\theta_1\approx Y/(Y+K_{d1})\approx Y/K_{d1}\ll1$, and the
    total target captured is $S\theta_1\approx SY/K_{d1}$, which is small
    because both $S\ll K_{d1}$ and $Y\ll K_{d1}$. Little DNA is recovered,
    so libraries have low complexity and are sensitive to contamination.
    \item \emph{Off-target contamination when $Y$ is raised.} Increasing $Y$
    to improve yield raises $\bar\rho$ from $r$ toward $Y/K_{d2}$ and
    eventually 1. Among endogenous reads, the target-to-off-target ratio is
    $A:\bar\rho C$, so the fraction of endogenous reads that are on target,
    $$
    \frac{A}{A+\bar\rho C},
    $$
    falls as $Y$ grows. Once $Y\gg K_{d1}$, target capture is saturated
    ($\theta_1\approx1$), so further increases in $Y$ only add off-target
    background. Because $\bar\rho$ is the same in both conditions, this
    background does not bias spike-in normalization at target loci (under
    the assumption that target and off-target loci do not overlap), but it
    lowers the signal-to-noise ratio, especially in the treatment
    condition, where the on-target fraction $A'/(A'+\bar\rho C)$ is
    smaller because $A'<A$.
\end{itemize}
